\documentclass{article}
\usepackage{graphicx}
\usepackage{float}
\usepackage{caption}
\usepackage{subcaption}
\usepackage{hyperref}
\usepackage{url}
\usepackage[export]{adjustbox} % for max width/height options
\usepackage{placeins} % for \FloatBarrier
\setlength{\parindent}{0pt}
\usepackage{verbatim}
\usepackage{amsmath}
\title{Rabbit Alert Modeling Project}
\author{Sonit Ambashta}
\date{\today}
\usepackage[a4paper,top=1cm, bottom=1cm, left=2cm, right=2cm]{geometry}

\captionsetup{justification=centering}

\begin{document}
\maketitle

\section{Variables}
\begin{itemize}
    \item $t$: time (in steps, can adapt to days/minutes/seconds as needed)
    \item $A(t)$: number of alerted rabbits at time $t$
    \item $k$: alert propagation rate, measuring how quickly alerted rabbits cause additional rabbits to become alerted
    \item $\frac{dA}{dt}$: instantaneous rate of change of the alerted rabbit population.
    \item $P$: maximum number of alerted rabbits (capacity of the habitat)
\end{itemize}

\section{Assumptions}
\begin{itemize}
\item The rabbit population contains at least one rabbit. 
\item One or more rabbits may initially detect the predator. 
\item Alert propagation occurs through interactions between rabbits and is modeled as proportional to the number of rabbits that are already alerted. 
\item The alert propagation rate $k$ is treated as a user-defined parameter because an empirically validated value is not currently available. 
\item All rabbits are equally capable of receiving and transmitting warning signals. 
\item Once alerted, a rabbit remains alerted for the duration of the simulation. 
\item Environmental effects such as terrain, vegetation density, weather, and visibility are ignored. 
\item Predator behavior and movement are not modeled. 
\item Time is treated as a continuous variable. 
\item The rabbit population is assumed to be well-mixed, meaning each alerted rabbit has an equal opportunity to influence any other rabbit.
\item The model is intended for wild rabbit populations and does not account for behaviors associated with domesticated rabbits.
\item Differences in alert behavior among rabbit species are not modeled. All wild rabbits are treated as having identical alert-propagation characteristics.
\end{itemize}

\section{Differential Equation and Approach}
\[ \frac{dA}{dt} = kA \]

Note: this is a separable differential equation
\begin{itemize}
    \item multiply $dt$ on both sides
    \[ dA = kA dt \]
    \item divide $A$ on both sides
    \[ \frac{1}{A}dA = k dt \]
    \item integrate both sides (don't forget the constant of integration)
    \[ \int\frac{1}{A}dA = \int k dt \]
    \[ ln |A| = kt + C \]
    \item isolate A
    \[ A = e^{kt + C}\]
    \item use exponent rules to separate the C 
    \[ A = e^C e^{kt} \] 
    \item let $e^C$ be an arbitrary constant $C_1$ 
    \[ A = C_1 e^{kt} \] 
    \item apply the initial condition $A(0)=A_0$ 
    \[ A_0 = C_1 e^{k(0)} \] 
    \item simplify since $e^0 = 1$ 
    \[ A_0 = C_1 \] 
    \item substitute $C_1 = A_0$ back into the solution 
    \[ A = A_0 e^{kt} \]
\end{itemize}

\subsection{Example Model} Suppose a user provides 
\[ k = 0.7 \] and \[ A_0 = 2. \] 
The differential equation becomes 
\[ \frac{dA}{dt} = 0.7A. \] 
Using the solution derived previously, 
\[ A(t)=A_0e^{kt}, \] 
we substitute the values of $A_0$ and $k$: 
\[ A(t)=2e^{0.7t}. \] 
This function models the number of alerted rabbits 
over time assuming that each alerted rabbit contributes 
proportionally to the spread of awareness throughout the population. 
The corresponding Python code is: 
\begin{verbatim} 
>>> equation = form_differential_equation(0.7)
>>> solution = solve_equation(equation, 2)
>>> print(solution)
\end{verbatim} 
Output: 
\begin{verbatim} 
Eq(A(t), 2*exp(0.7*t))
\end{verbatim} 
Some example evaluations are: 
\[ A(0)=2e^0=2 \] \[ A(1)=2e^{0.7}\approx 4.03 \] \[ A(2)=2e^{1.4}\approx 8.11 \] \[ A(3)=2e^{2.1}\approx 16.33 \] These values illustrate the exponential nature of the model, where the number of alerted rabbits grows increasingly rapidly over time. This behavior motivates the introduction of a population cap in later versions of the model.
\subsection{Notes}
\begin{itemize}
    \item This equation is for any arbitrary k that will be plugged in (user adjustable), however growth cannot be demonstrated for $k = 0$ since
that would signify no change.
    \item Right now, there is not a supporting rate at which rabbits get alerted, so this serves as a user-adjustable cli interface that allows adding a rate
    \item A certain number of rabbits can be alerted at a moment of time (especially when they are together), so the initial value setup can also be provided by the user
    \item In this model, it is also assumed that there is no upper bound on the number of rabbits that can be alerted at any given time.
\end{itemize}

\section{Major Issue with this Approach and Proposed Solution} 
The original model assumed an unlimited rabbit population, producing unbounded exponential growth. 
This assumption was later relaxed through the introduction of a carrying-capacity parameter $P$, 
resulting in a capped exponential model. The idea is to allow exponential growth until the number 
of alerted rabbits reaches the maximum number of rabbits in the habitat. To determine when this occurs, 
solve for the point of intersection between the exponential solution and the carrying capacity.
 \[ A(t)=A_0e^{kt} \] \[ P=A_0e^{kt} \] 
 Divide both sides by $A_0$: 
 \[ \frac{P}{A_0}=e^{kt} \] 
 Take the natural logarithm of both sides:
\[ \ln\left(\frac{P}{A_0}\right)=kt \] 
 Solve for $t$: 
 \[ t^*=\frac{\ln(P/A_0)}{k} \] 
 where $t^*$ represents the time at which the alerted population reaches the habitat capacity. 
 Once this point is reached, there are no additional rabbits available to become alerted. Therefore, 
 the rate of change becomes zero. This gives the piecewise differential equation 
 \[ \frac{dA}{dt} = \begin{cases} kA, & A < P, \\ 0, & A \ge P. \end{cases} \] 
 The corresponding piecewise solution is 
 \[ A(t) = \begin{cases} A_0e^{kt}, & t < t^*, \\ P, & t \ge t^*. \end{cases} \] 
 where \[ t^*=\frac{\ln(P/A_0)}{k}. \]

\section{Inspiration and Rationale}
\begin{itemize}
    \item This project began as an opportunity to apply differential equations outside of a classroom setting and 
explore how mathematical models can be used to represent real-world phenomena programmatically. Rabbits 
were chosen as the motivating domain because they are an important component of many ecosystems and 
serve as prey for a wide range of predators. Their survival often depends on the rapid detection and 
communication of threats, making predator-awareness propagation a natural phenomenon to explore through 
differential equations. The initial goal was not to produce a biologically validated model. Rather, 
the objective was to practice the process of mathematical modeling: defining assumptions, constructing 
a differential equation, solving the equation analytically, implementing the solution in Python, and 
evaluating the realism and limitations of the resulting model. 
    \item As development progressed, the project evolved beyond simply solving an exponential differential equation. 
The model was extended to include a finite habitat capacity after it became clear that unrestricted exponential 
growth implied an unrealistic infinite rabbit population. This led to the introduction of a capped exponential 
model and further discussion regarding alternative modeling approaches and their underlying assumptions. 
    \item Although a validated behavioral propagation rate is not currently available, the framework is designed as a 
command-line utility for exploring different parameter combinations. The project may serve as an educational tool 
for learning differential equations and mathematical modeling, while also providing a foundation for future 
investigations into rabbit behavior, ecosystem dynamics, and animal communication.
\end{itemize}

\end{document}